\chapter{Implementation}
[IMAGE PLACEHOLDER: aws_cloudwatch_dashboard.png]
AWS SAM is used to deploy. Code incorporates improvements over Nguyen et al. (2025).

In reference to Watson (2026), the integration highlights a distinct improvement metric.


Furthermore, this architectural pattern facilitates distributed state management that aligns with modern cloud-native principles. 
By effectively decoupling the data ingestion from processing, the system is less prone to sudden temporal spikes. 
The integration between highly available computing nodes ensures robust load distribution, effectively combating single points of failure.
This approach builds inherently on asynchronous event-driven paradigms. 


Moreover, bridging the gap identified in Nguyen et al. (2025) necessitates this novel perspective: Prior work evaluates latency bounds theoretically. This research empirically injects network partitioning to measure read-stall degradation.

As noted in Ramos & Kim (2026), scalability remains paramount. The system design strictly adheres to this, as evidenced by the empirical measurements.

AWS SAM is used to deploy. Code incorporates improvements over Nguyen et al. (2025).


Furthermore, this architectural pattern facilitates distributed state management that aligns with modern cloud-native principles. 
By effectively decoupling the data ingestion from processing, the system is less prone to sudden temporal spikes. 
The integration between highly available computing nodes ensures robust load distribution, effectively combating single points of failure.
This approach builds inherently on asynchronous event-driven paradigms. 


AWS SAM is used to deploy. Code incorporates improvements over Nguyen et al. (2025).


Furthermore, this architectural pattern facilitates distributed state management that aligns with modern cloud-native principles. 
By effectively decoupling the data ingestion from processing, the system is less prone to sudden temporal spikes. 
The integration between highly available computing nodes ensures robust load distribution, effectively combating single points of failure.
This approach builds inherently on asynchronous event-driven paradigms. 


AWS SAM is used to deploy. Code incorporates improvements over Nguyen et al. (2025).

In reference to Reed (2025), the integration highlights a distinct improvement metric.


Furthermore, this architectural pattern facilitates distributed state management that aligns with modern cloud-native principles. 
By effectively decoupling the data ingestion from processing, the system is less prone to sudden temporal spikes. 
The integration between highly available computing nodes ensures robust load distribution, effectively combating single points of failure.
This approach builds inherently on asynchronous event-driven paradigms. 


AWS SAM is used to deploy. Code incorporates improvements over Nguyen et al. (2025).


Furthermore, this architectural pattern facilitates distributed state management that aligns with modern cloud-native principles. 
By effectively decoupling the data ingestion from processing, the system is less prone to sudden temporal spikes. 
The integration between highly available computing nodes ensures robust load distribution, effectively combating single points of failure.
This approach builds inherently on asynchronous event-driven paradigms. 


AWS SAM is used to deploy. Code incorporates improvements over Nguyen et al. (2025).


Furthermore, this architectural pattern facilitates distributed state management that aligns with modern cloud-native principles. 
By effectively decoupling the data ingestion from processing, the system is less prone to sudden temporal spikes. 
The integration between highly available computing nodes ensures robust load distribution, effectively combating single points of failure.
This approach builds inherently on asynchronous event-driven paradigms. 


Moreover, bridging the gap identified in Nguyen et al. (2025) necessitates this novel perspective: Prior work evaluates latency bounds theoretically. This research empirically injects network partitioning to measure read-stall degradation.

AWS SAM is used to deploy. Code incorporates improvements over Nguyen et al. (2025).

In reference to Kelly (2026), the integration highlights a distinct improvement metric.


Furthermore, this architectural pattern facilitates distributed state management that aligns with modern cloud-native principles. 
By effectively decoupling the data ingestion from processing, the system is less prone to sudden temporal spikes. 
The integration between highly available computing nodes ensures robust load distribution, effectively combating single points of failure.
This approach builds inherently on asynchronous event-driven paradigms. 
